\section{Robustness and Extensions}
\label{sec:robustness}

We subject our findings to extensive robustness checks organized in six
groups. Section~\ref{sec:robustness}.1 varies the FL definition
threshold, win-rate cutoff, and temporal window.
Section~\ref{sec:robustness}.2 checks sample restrictions and
specification variants (unrestricted sample, homogeneous subsamples,
tighter controls, matching, clustering).
Section~\ref{sec:robustness}.3 runs sensitivity analysis on unobservable
confounders. Section~\ref{sec:robustness}.4 reports the staggered DiD
(complementary, underpowered). Section~\ref{sec:robustness}.5 computes
welfare bounds under OLS, IV, and cross-fit estimates.
Section~\ref{sec:robustness}.6 examines procurement oversight and
compliance channels.

\subsection{FL Definition Robustness}

\paragraph{IQR threshold sensitivity.}
We vary the IQR multiplier from 1.0 to 3.0, which changes both the
threshold and the number of FL firms (from 3,442 at 1.0$\times$ to 1,456
at 3.0$\times$). The price coefficient increases monotonically with the
multiplier, consistent with a stricter threshold selecting more extreme
cover bidders.
Figure~\ref{fig:threshold_stability} and Table~\ref{tab:threshold} in
Appendix~\ref{app:robustness} report details.

\paragraph{Low-win-rate FL variants.}
Instead of requiring a zero win rate, we define FL using win rate
$< 1\%$ and $< 2\%$. These variants capture firms that win trivially
(1--2 wins in hundreds of participations) but otherwise behave as cover
bidders. Table~\ref{tab:fl_lowwinrate} in Appendix~\ref{app:robustness}
shows that coefficients are stable across these definitions.

\paragraph{Cross-fitting.}
We define FL firms using only odd years (2009, 2011, \ldots, 2019) and
estimate regressions on even years, and vice versa. This breaks any
mechanical link between FL classification and outcomes.
Table~\ref{tab:fl_crossfit} in Appendix~\ref{app:robustness} shows that
the odd-year-trained fold yields a significant coefficient of 0.053 % R1-FIX: deeper engagement with cross-fit attenuation
($p < 0.05$); the even-year-trained fold yields 0.019, which is
positive but not individually significant. The cross-fold average of
0.036 is attenuated by 44\% relative to the full-sample OLS (0.064).

Three factors contribute to the attenuation and fold asymmetry. First,
the even-year fold identifies fewer FL firms (1,885 vs.\ 2,153 in
the odd fold) because even years contain fewer observations for
threshold computation, producing a lower threshold (12 vs.\ 13.5)
that admits more marginal firms---diluting the FL signal. Second,
defining FL on a temporal subsample introduces \emph{classification
noise}: firms near the threshold may switch FL status across training
folds, generating measurement-error attenuation in the cross-fit
estimate relative to the full-sample estimate. Third, any mechanical
correlation between FL classification and contemporaneous prices
(e.g., FL firms selecting into the same tenders they are classified
on) is broken by the temporal split, so the cross-fit average isolates
the \emph{out-of-sample} predictive content of FL status.

We interpret the cross-fit average (0.036, or 3.6\%) as a conservative
lower bound on the true FL--price association. The convergence of the
cross-fit (3.6\%), structural (4.3\%), and full-sample OLS (6.4\%)
estimates---all within a factor of two---suggests that the FL screen
captures a genuine price signal, with the variation across estimators
reflecting different degrees of measurement-error attenuation rather
than a fundamentally spurious association.

\paragraph{Temporal FL definition.}
Instead of classifying FL firms over the full 2009--2019 period, we use
rolling 3-year windows: a firm is FL in year $t$ if it is an
always-loser with above-threshold participation in years $[t-2, t]$.
Table~\ref{tab:fl_temporal} in Appendix~\ref{app:robustness} shows
that time-varying FL status produces a similar coefficient of 0.070
($p < 0.01$, $N = 1{,}255{,}998$), close to the baseline 0.064.

\subsection{Sample and Specification Robustness}

\paragraph{Unrestricted sample.}
Our baseline restricts to item types with at least one FL-present
tender. Removing this restriction yields a coefficient of 0.063 ($N =
2{,}875{,}653$), virtually identical to the baseline 0.064
(Table~\ref{tab:unrestricted_sample} in Appendix~\ref{app:robustness}),
suggesting that restriction (iii) does not bias the results.

\paragraph{Homogeneous sub-samples.}
Identifying \emph{ex ante} homogeneous goods from item descriptions is
infeasible given BEC's coding structure. Instead, we use a data-driven
approach: for each item group (first two digits of the item code), we
compute the within-group coefficient of variation (CV) of log negotiated
prices. Groups with low CV contain items with similar prices across
tenders---a proxy for product homogeneity. We split item groups at the
median CV and run FL regressions on each subsample.
Table~\ref{tab:homogeneous_cv} in Appendix~\ref{app:robustness} reports
the results. If the FL effect is driven by quality differences rather
than cover bidding, we would expect it to vanish in low-CV (homogeneous)
groups where quality-based selection stories are implausible. Item fixed
effects already absorb much cross-item heterogeneity, so the CV split
provides an additional dimension of variation.

\paragraph{Tighter controls.}
We add controls beyond item/year/PBU FE.
Table~\ref{tab:tighter_controls} in Appendix~\ref{app:robustness}
reports four specifications. Adding the number of genuine bidders
increases the coefficient to 0.094 (column~2), consistent with FL
presence being positively selected into competitive tenders.
Item $\times$ year FE (column~4) yield 0.074, close to the baseline.
Column~(3) adds log reference price, which absorbs most cross-tender
price variation ($R^2$ rises from 0.886 to 0.973); the FL coefficient
becomes $-0.031$. This sign reversal is consistent with a bad-control
problem: if cover bidding inflates reference prices---e.g., through
prior-tender manipulation---then conditioning on reference price removes
the treatment effect and may induce collider bias. We therefore treat
column~(3) as a mechanical decomposition of the price level rather than
a robustness check.

\paragraph{Matching estimators.} % V7-T2: brief reference, detail in Section 7
CEM ($\hat{\beta} = 0.077$) and IPW ($\hat{\beta} = 0.055$) are
reported as primary estimates in Section~\ref{sec:results_ols};
balance diagnostics are in Table~\ref{tab:matching} in
Appendix~\ref{app:robustness}.

\paragraph{Alternative clustering.}
Standard errors clustered at the item level (baseline: 0.022), PBU level
(0.014), and two-way item + PBU level (0.024) all leave inference
unchanged, with the coefficient significant at the 1\% level under all
clustering schemes.

\subsection{Sensitivity to Unobservables}

\cites{cinelli2020making} sensitivity framework quantifies how strong
omitted variable bias would need to be to explain away the FL
coefficient. The robustness value $RV_{q=1} = 17.5\%$ means that an
unobserved confounder would need to explain at least 17.5\% of the
residual variation in both FL presence and log prices to reduce the
coefficient to zero---a stringent requirement given the rich fixed
effects structure.

\subsection{Staggered Difference-in-Differences}

We report the staggered DiD as a complementary exercise only. Using
\citet{callaway2021difference}, we construct a market-level panel (item
$\times$ PBU $\times$ year) with 144,168 market-year observations across
19,777 markets (1,511 treated, 18,266 never-treated). The C\&S estimator
yields an overall ATT of 0.014 (SE~=~0.039) for log prices---positive
but insignificant. The minimum detectable effect is
$1.96 \times 0.039 = 0.076$, exceeding the OLS estimate of 0.064.
Given limited treatment timing variation, the design is underpowered for
detecting effects of the magnitude we estimate via OLS and IV, and we do
not rely on it for identification.

\cites{rambachan2023more} sensitivity analysis, sharp-entry subsamples,
and event-study plots are reported in Appendix~\ref{app:did}
(Table~\ref{tab:did_revised} and Figure~\ref{fig:event_study}).

\subsection{Welfare Analysis}
\label{sec:welfare}

We compute welfare loss bounds from cover bidding using the formula
$\text{Welfare loss} = (\exp(\hat{\beta}) - 1) \times
\sum_{i: \text{FL}_i = 1} p_i$.

The OLS coefficient provides a \emph{lower bound} (if measurement error
attenuates $\hat{\beta}$). The IV coefficient provides an
\emph{upper bound} (correcting for attenuation). The concentrated-market
network coefficient provides a targeted estimate for the subset of FL
firms most likely to be genuine cover bidders.

Table~\ref{tab:welfare_bounds} in Appendix~\ref{app:extensions} reports
the welfare bounds. The OLS-based estimate is R\$734 million (95\% CI:
R\$243M--R\$1,246M), providing a lower bound under measurement-error
attenuation. The cross-fit average coefficient (0.036) provides an
intermediate bound: by defining FL using odd years and estimating on even
years (and vice versa), this breaks any mechanical link between
classification and outcomes, yielding a more conservative estimate. The
IV-based estimate is R\$2.4 billion, reflecting the larger instrumented
coefficient. As a share of total procurement spending (R\$30.8 billion), the OLS
welfare loss represents approximately 2.4\% of total BEC procurement
value over 2009--2019, the cross-fit estimate 1.3\%, and the IV estimate
7.8\%.

These figures should be interpreted as bounds: the OLS estimate is
conservative (attenuated by binary treatment misclassification), the
cross-fit estimate is intermediate, and the IV estimate captures the
Local Average Treatment Effect (LATE) for complier tenders---those whose
FL participation status is shifted by supply-side availability. The LATE
need not equal the ATE: if complier tenders are disproportionately
affected by cover bidding (e.g., in markets where FL supply variation is
most consequential), the IV-based welfare bound may overstate the average
effect across all FL-present tenders.

\subsection{Oversight Heterogeneity}

Table~\ref{tab:regime_oversight} in Appendix~\ref{app:robustness}
reports the FL price coefficient by PBU size quartile. The coefficient
is largest for the smallest PBUs (Q1: 0.214, imprecisely estimated) and
smallest for the largest PBUs (Q4: 0.017, not individually significant),
broadly consistent with the conceptual framework's prediction that
detection probability $\theta$ constrains cover bidding. Only Q3 (0.066)
is individually significant ($p < 0.05$). The pattern across quartiles
is noisy, reflecting limited variation in FL presence within small PBUs.

\begin{figure}[htbp]
\centering
\includegraphics[width=0.85\textwidth]{fig_07_threshold_stability.pdf}
\caption{FL price coefficient across IQR multiplier thresholds.
Error bars show 95\% confidence intervals.}
\label{fig:threshold_stability}
\end{figure}
